% Chunk: §7.5 Transition to Interaction Picture: Examples (through Eq. 7.5.48)
% Source: Weinberg QFT Vol. I, printed pp. 318--324 (physical PDF pp. 341--347)
% Status: transcribed; source-reviewed; compile-clean; render-reviewed
% Figures: none. Uncertainties: none.

\subsection{Transition to Interaction Picture: Examples}

At the end of Section 7.2 we showed how to use the Lagrangian of a
simple scalar field theory to derive the structure of the interaction and
the free fields it contains in the interaction picture. We will now turn to
somewhat more complicated and revealing examples.

\subsubsection*{Scalar Field, Derivative Coupling}

First let's consider a neutral scalar field, but now with derivative coupling.
We take the Lagrangian as
\begin{equation}
\mathcal L=-\frac12\partial_\mu\Phi\partial^\mu\Phi
-\frac12m^2\Phi^2-J^\mu\partial_\mu\Phi-\mathcal H_I(\Phi),
\tag{7.5.1}\label{eq:7.5.1}
\end{equation}
where $J^\mu$ is either a c-number external current (unrelated to currents
$J^\mu$ introduced earlier), or a functional of various fields other than
$\Phi$ (in which case terms involving these other fields need to be added to
\eqref{eq:7.5.1}). The canonical conjugate to $\Phi$ is now
\begin{equation}
\Pi=\frac{\partial\mathcal L}{\partial\dot\Phi}=\dot\Phi-J^0
\tag{7.5.2}\label{eq:7.5.2}
\end{equation}
and the Hamiltonian is
\begin{align*}
H&=\int \dd^3x\,[\Pi\dot\Phi-\mathcal L]\notag\\
&=\int \dd^3x\Bigl[\Pi(\Pi+J^0)+\frac12(\boldsymbol\nabla\Phi)^2
-\frac12(\Pi+J^0)^2+\frac12m^2\Phi^2\notag\\
&\hspace{7.8em}{}+\mathbf J\cdot\boldsymbol\nabla\Phi
+J^0(\Pi+J^0)+\mathcal H_I(\Phi)\Bigr].
\end{align*}
Collecting terms, we can write this as
\begin{equation}
H=H_0+H_{\mathrm{int}},
\tag{7.5.3}\label{eq:7.5.3}
\end{equation}
\begin{equation}
H_0=\int \dd^3x\left[\frac12\Pi^2+\frac12(\boldsymbol\nabla\Phi)^2
+\frac12m^2\Phi^2\right],
\tag{7.5.4}\label{eq:7.5.4}
\end{equation}
\begin{equation}
H_{\mathrm{int}}=\int \dd^3x\left[\Pi J^0+\mathbf J\cdot\boldsymbol\nabla\Phi
+\frac12(J^0)^2+\mathcal H_I(\Phi)\right].
\tag{7.5.5}\label{eq:7.5.5}
\end{equation}
As explained in Section 7.2, we can pass to the interaction picture by
simply replacing $\Pi$ and $\Phi$ with $\pi$ and $\phi$ (and likewise for any
fields in the current $J^\mu$, though we will not bother to indicate this
explicitly):
\begin{equation}
H_0=\int \dd^3x\left[\frac12\pi^2(\mathbf x,t)
+\frac12\bigl(\boldsymbol\nabla\phi(\mathbf x,t)\bigr)^2
+\frac12m^2\phi^2(\mathbf x,t)\right],
\tag{7.5.6}\label{eq:7.5.6}
\end{equation}
\begin{align}
H_I(t)=\int \dd^3x\Bigl[&\pi(\mathbf x,t)J^0(\mathbf x,t)
+\mathbf J(\mathbf x,t)\cdot\boldsymbol\nabla\phi(\mathbf x,t)\notag\\
&+\frac12\bigl[J^0(\mathbf x,t)\bigr]^2
+\mathcal H_I\bigl(\phi(\mathbf x,t)\bigr)\Bigr].
\tag{7.5.7}\label{eq:7.5.7}
\end{align}
The free-particle Hamiltonian is just the same as Eq.~\eqref{eq:7.2.25}, and leads
as in Section 7.2 to Eqs.~\eqref{eq:7.2.26}--\eqref{eq:7.2.35}. Indeed, whatever the total
Hamiltonian may be, we \emph{must} take \eqref{eq:7.5.6} as the part we split off and
call the free-particle part, with the remainder called the interaction,
because as we have seen it is this form of the free-particle Hamiltonian that
leads to the correct expansion \eqref{eq:7.2.29} of the scalar field in terms of
creation and annihilation operators that satisfy the commutation relations
\eqref{eq:7.2.34}, \eqref{eq:7.2.35}. The last step is to replace $\pi$ in the interaction
Hamiltonian with its value $\dot\phi$ in the interaction picture (\emph{not}
its value $\dot\Phi-J^0$ in the Heisenberg picture):
\begin{equation}
H_I(t)=\int \dd^3x\left[J^\mu(\mathbf x,t)\partial_\mu\phi(\mathbf x,t)
+\frac12\bigl[J^0(\mathbf x,t)\bigr]^2
+\mathcal H_I\bigl(\phi(\mathbf x,t)\bigr)\right].
\tag{7.5.8}\label{eq:7.5.8}
\end{equation}
The extra non-invariant term in Eq.~\eqref{eq:7.5.8} is just what we saw in Section
6.2 is needed to cancel a non-invariant term in the propagator of
$\partial_0\phi$.

\subsubsection*{Vector Field, Spin One}

Similar results are obtained in the canonical quantization of the vector
field $V_\mu$ for a particle of spin one. Let's here keep an open mind, and
write the Lagrangian density in a fairly general form
\begin{equation}
\mathcal L=-\frac12\alpha\,\partial_\mu V_\nu\partial^\mu V^\nu
-\frac12\beta\,\partial_\mu V_\nu\partial^\nu V^\mu
-\frac12m^2V_\mu V^\mu-J_\mu V^\mu,
\tag{7.5.9}\label{eq:7.5.9}
\end{equation}
where $\alpha$, $\beta$, and $m^2$ are so far arbitrary constants, and
$J_\mu$ is either a c-number external current, or an operator depending on
fields other than $V^\mu$, in which case additional terms involving these
fields must be added to $\mathcal L$. The Euler--Lagrange field equations for
$V_\mu$ read
\[
\alpha\Box V^\nu+\beta\partial^\nu(\partial_\mu V^\mu)+m^2V^\nu=-J^\nu.
\]
Taking the divergence gives
\begin{equation}
(\alpha+\beta)\Box\partial_\lambda V^\lambda
+m^2\partial_\lambda V^\lambda=-\partial_\lambda J^\lambda.
\tag{7.5.10}\label{eq:7.5.10}
\end{equation}
This is the equation for an ordinary scalar field with mass
$m^2/(\alpha+\beta)$ and source $\partial_\lambda J^\lambda/(\alpha+\beta)$.
We want to describe a theory containing only particles of spin one, not spin
zero, so to avoid the appearance of $\partial_\lambda V^\lambda$ as an
independently propagating scalar field, we take $\alpha=-\beta$, in which case
$\partial_\lambda V^\lambda$ can be expressed in terms of an external current
or other fields, as $-\partial_\lambda J^\lambda/m^2$. The constant $\alpha$
can be absorbed in the definition of $V_\mu$, so we can take
$\alpha=-\beta=1$, and therefore
\begin{equation}
\mathcal L=-\frac14F_{\mu\nu}F^{\mu\nu}
-\frac12m^2V_\mu V^\mu-J_\mu V^\mu,
\tag{7.5.11}\label{eq:7.5.11}
\end{equation}
where
\begin{equation}
F_{\mu\nu}=\partial_\mu V_\nu-\partial_\nu V_\mu.
\tag{7.5.12}\label{eq:7.5.12}
\end{equation}
The derivative of the Lagrangian with respect to the time-derivative of the
vector field is
\begin{equation}
\frac{\partial\mathcal L}{\partial\dot V^\mu}=-F^{0\mu}.
\tag{7.5.13}\label{eq:7.5.13}
\end{equation}
This is non-vanishing for $\mu$ a spatial index $i$, so the $V^i$ are
canonical fields, with conjugates
\begin{equation}
\Pi^i=F^{i0}=\dot V^i+\partial_iV^0.
\tag{7.5.14}\label{eq:7.5.14}
\end{equation}
On the other hand $F^{00}=0$, so $\dot V^0$ does not appear in the
Lagrangian, and $V^0$ is therefore an auxiliary field. This causes no serious
difficulty: the fact that $\partial\mathcal L/\partial\dot V^0$ vanishes
means that the field equation for $V^0$ involves no second time-derivatives,
and can therefore be used as a constraint that eliminates a field variable.
Specifically, the Euler--Lagrange equation for $\nu=0$ is
\begin{equation}
\partial_iF^{i0}=m^2V^0+J^0
\tag{7.5.15}\label{eq:7.5.15}
\end{equation}
or using Eq.~\eqref{eq:7.5.14}
\begin{equation}
V^0=\frac1{m^2}(\boldsymbol\nabla\cdot\boldsymbol\Pi-J^0).
\tag{7.5.16}\label{eq:7.5.16}
\end{equation}

Now let us calculate the Hamiltonian
$H=\int \dd^3x\,(\boldsymbol\Pi\cdot\dot{\mathbf V}-\mathcal L)$ for this
theory. Eq.~\eqref{eq:7.5.14} allows us to write $\dot{\mathbf V}$ in terms of
$\boldsymbol\Pi$ and $J^0$:
\[
\dot{\mathbf V}=-\boldsymbol\nabla V^0+\boldsymbol\Pi
=\boldsymbol\Pi-\frac1{m^2}\boldsymbol\nabla
(\boldsymbol\nabla\cdot\boldsymbol\Pi-J^0),
\]
so
\begin{align*}
H=\int \dd^3x\Bigl[&\boldsymbol\Pi^2
+m^{-2}(\boldsymbol\nabla\cdot\boldsymbol\Pi)
(\boldsymbol\nabla\cdot\boldsymbol\Pi-J^0)
-\frac12\boldsymbol\Pi^2+\frac12(\boldsymbol\nabla\times\mathbf V)^2
+\frac12m^2\mathbf V^2\notag\\
&-\frac12m^{-2}(\boldsymbol\nabla\cdot\boldsymbol\Pi-J^0)^2
+\mathbf J\cdot\mathbf V
-m^{-2}J^0(\boldsymbol\nabla\cdot\boldsymbol\Pi-J^0)\Bigr].
\end{align*}
Again, we split this up into a free-particle term $H_0$ and interaction
$H_{\mathrm{int}}$:
\begin{equation}
H=H_0+H_{\mathrm{int}},
\tag{7.5.17}\label{eq:7.5.17}
\end{equation}
and pass to the interaction picture by replacing the Heisenberg-picture
quantities $\mathbf V$ and $\boldsymbol\Pi$ with their interaction-picture
counterparts $\mathbf v$ and $\boldsymbol\pi$ (and, though not shown
explicitly, likewise for whatever fields and conjugates are present in
$J^\mu$):
\begin{equation}
H_0=\int \dd^3x\left[\frac12\boldsymbol\pi^2
+\frac1{2m^2}(\boldsymbol\nabla\cdot\boldsymbol\pi)^2
+\frac12(\boldsymbol\nabla\times\mathbf v)^2
+\frac{m^2}{2}\mathbf v^2\right],
\tag{7.5.18}\label{eq:7.5.18}
\end{equation}
\begin{equation}
H_I(t)=\int \dd^3x\left[\mathbf J\cdot\mathbf v
-m^{-2}J^0\boldsymbol\nabla\cdot\boldsymbol\pi
+\frac1{2m^2}(J^0)^2\right].
\tag{7.5.19}\label{eq:7.5.19}
\end{equation}
The relation between $\boldsymbol\pi$ and $\mathbf v$ is then
\begin{equation}
\dot{\mathbf v}=\frac{\delta H_0(\mathbf v,\boldsymbol\pi)}
{\delta\boldsymbol\pi}
=\boldsymbol\pi-m^{-2}\boldsymbol\nabla
(\boldsymbol\nabla\cdot\boldsymbol\pi)
\tag{7.5.20}\label{eq:7.5.20}
\end{equation}
and the `field equation' is
\begin{equation}
\dot{\boldsymbol\pi}=-\frac{\delta H_0(\mathbf v,\boldsymbol\pi)}
{\delta\mathbf v}
=+\nabla^2\mathbf v-\boldsymbol\nabla(\boldsymbol\nabla\cdot\mathbf v)
-m^2\mathbf v.
\tag{7.5.21}\label{eq:7.5.21}
\end{equation}
Since $V^0$ is not an independent field variable, it is not related by a
similarity transformation to any interaction-picture object $v^0$. Instead,
we can \emph{invent} a quantity
\begin{equation}
v^0\equiv m^{-2}\boldsymbol\nabla\cdot\boldsymbol\pi.
\tag{7.5.22}\label{eq:7.5.22}
\end{equation}
Eq.~\eqref{eq:7.5.20} then allows us to write $\boldsymbol\pi$ as
\begin{equation}
\boldsymbol\pi=\dot{\mathbf v}+\boldsymbol\nabla v^0.
\tag{7.5.23}\label{eq:7.5.23}
\end{equation}
Inserting this in Eqs.~\eqref{eq:7.5.22} and \eqref{eq:7.5.21} gives our field equations in the
form
\begin{align*}
\nabla^2v^0+\boldsymbol\nabla\cdot\dot{\mathbf v}-m^2v^0&=0,\\
\nabla^2\mathbf v-\boldsymbol\nabla(\boldsymbol\nabla\cdot\mathbf v)
-\ddot{\mathbf v}-\boldsymbol\nabla\dot v^0-m^2\mathbf v&=0.
\end{align*}
These can be combined in the covariant form
\begin{equation}
\Box v^\mu-\partial^\mu\partial_\nu v^\nu-m^2v^\mu=0.
\tag{7.5.24}\label{eq:7.5.24}
\end{equation}
Taking the divergence gives
\begin{equation}
\partial_\mu v^\mu=0
\tag{7.5.25}\label{eq:7.5.25}
\end{equation}
and hence
\begin{equation}
(\Box-m^2)v^\mu=0.
\tag{7.5.26}\label{eq:7.5.26}
\end{equation}
A real vector field satisfying Eqs.~\eqref{eq:7.5.25} and \eqref{eq:7.5.26} can be expressed as
a Fourier transform
\begin{align}
v^\mu(x)=(2\pi)^{-3/2}\sum_\sigma\int \dd^3p\,(2p^0)^{-1/2}
\Bigl\{&e^\mu(\mathbf p,\sigma)a_{\mathbf p,\sigma}e^{\ii p\cdot x}\notag\\
&+e^{\mu*}(\mathbf p,\sigma)a^\dagger_{\mathbf p,\sigma}e^{-\ii p\cdot x}\Bigr\},
\tag{7.5.27}\label{eq:7.5.27}
\end{align}
where $p^0=\sqrt{\mathbf p^2+m^2}$; the $e^\mu(\mathbf p,\sigma)$ for
$\sigma=+1,0,-1$ are three independent vectors satisfying
\begin{equation}
p_\mu e^\mu(\mathbf p,\sigma)=0
\tag{7.5.28}\label{eq:7.5.28}
\end{equation}
and normalized so that
\begin{equation}
\sum_\sigma e^\mu(\mathbf p,\sigma)e^{\nu*}(\mathbf p,\sigma)
=\eta^{\mu\nu}+\frac{p^\mu p^\nu}{m^2};
\tag{7.5.29}\label{eq:7.5.29}
\end{equation}
and the $a_{\mathbf p,\sigma}$ are operator coefficients. It is straightforward
using Eqs.~\eqref{eq:7.5.23}, \eqref{eq:7.5.27}, and \eqref{eq:7.5.29} to calculate that $\mathbf v$ and
$\boldsymbol\pi$ satisfy the correct commutation relations
\begin{align}
[v^i(\mathbf x,t),\pi^j(\mathbf y,t)]&=\ii\delta_{ij}\delta^3(\mathbf x-\mathbf y),
\notag\\
[v^i(\mathbf x,t),v^j(\mathbf x,t)]
&=[\pi^i(\mathbf x,t),\pi^j(\mathbf x,t)]=0,
\tag{7.5.30}\label{eq:7.5.30}
\end{align}
provided that $a_{\mathbf p,\sigma}$ and $a^\dagger_{\mathbf p,\sigma}$ satisfy
the commutation relations
\begin{equation}
[a_{\mathbf p,\sigma},a^\dagger_{\mathbf p',\sigma'}]
=\delta^3(\mathbf p'-\mathbf p)\delta_{\sigma'\sigma},
\tag{7.5.31}\label{eq:7.5.31}
\end{equation}
\begin{equation}
[a_{\mathbf p,\sigma},a_{\mathbf p',\sigma'}]=0.
\tag{7.5.32}\label{eq:7.5.32}
\end{equation}

We already know that the vector field for a spin one particle must take the
form \eqref{eq:7.5.27}, so our derivation of these results serves to verify that Eq.
\eqref{eq:7.5.18} gives the correct free-particle Hamiltonian for a massive particle
of spin one. It is easy to check also that Eq.~\eqref{eq:7.5.18} may be written (up to
a constant term) in the standard form of a free-particle energy, as
$\sum_\sigma\int \dd^3p\,p^0a^\dagger_{\mathbf p,\sigma}a_{\mathbf p,\sigma}$.
Finally, using Eq.~\eqref{eq:7.5.22} in Eq.~\eqref{eq:7.5.19} yields the interaction in the
interaction picture
\begin{equation}
H_I(t)=\int \dd^3x\left[J_\mu v^\mu+\frac1{2m^2}(J^0)^2\right].
\tag{7.5.33}\label{eq:7.5.33}
\end{equation}
The extra non-invariant term in Eq.~\eqref{eq:7.5.33} is just what we found in Chapter
6 is needed to cancel a non-invariant term in the propagator of the vector
field.

\subsubsection*{Dirac field, Spin One Half}

For the Dirac field of a particle of spin $1/2$, we tentatively take the
Lagrangian as
\begin{equation}
\mathcal L=-\bar\Psi(\gamma^\mu\partial_\mu+m)\Psi
-\mathcal H_I(\bar\Psi,\Psi)
\tag{7.5.34}\label{eq:7.5.34}
\end{equation}
with $\mathcal H_I$ a real function of $\bar\Psi$ and $\Psi$. This is not real,
but the action is, because
\[
\bar\Psi\gamma^\mu\partial_\mu\Psi
-(\bar\Psi\gamma^\mu\partial_\mu\Psi)^\dagger
=\bar\Psi\gamma^\mu\partial_\mu\Psi
+(\partial_\mu\bar\Psi)\gamma^\mu\Psi
=\partial_\mu(\bar\Psi\gamma^\mu\Psi).
\]
Hence the field equations obtained by requiring the action to be stationary
with respect to $\bar\Psi$ are the adjoints of those obtained by requiring the
action to be stationary with respect to $\Psi$, as necessary if we are to
avoid having too many field equations. The canonical conjugate to $\Psi$ is
\begin{equation}
\Pi=\frac{\partial\mathcal L}{\partial\dot\Psi}=-\bar\Psi\gamma^0,
\tag{7.5.35}\label{eq:7.5.35}
\end{equation}
so we should not regard $\bar\Psi$ as a field like $\Psi$, but rather as
proportional to the canonical conjugate of $\Psi$. The Hamiltonian is
\[
H=\int \dd^3x\,[\Pi\dot\Psi-\mathcal L]
=\int \dd^3x\left[\Pi\gamma^0(\boldsymbol\gamma\cdot\boldsymbol\nabla+m)\Psi
+\mathcal H_I\right].
\]
We write this as
\begin{equation}
H=H_0+H_{\mathrm{int}},
\tag{7.5.36}\label{eq:7.5.36}
\end{equation}
where
\begin{equation}
H_0=\int \dd^3x\,\Pi\gamma^0(\boldsymbol\gamma\cdot\boldsymbol\nabla+m)\Psi,
\tag{7.5.37}\label{eq:7.5.37}
\end{equation}
\begin{equation}
H_{\mathrm{int}}=\int \dd^3x\,\mathcal H_I(\bar\Psi,\Psi).
\tag{7.5.38}\label{eq:7.5.38}
\end{equation}

We now pass to the interaction picture. Since Eq.~\eqref{eq:7.5.35} does not involve
the time, the similarity transformation \eqref{eq:7.1.28}, \eqref{eq:7.1.29} yields immediately
\begin{equation}
\pi=-\bar\psi\gamma^0.
\tag{7.5.39}\label{eq:7.5.39}
\end{equation}
Likewise, $H_0$ and $H_I(t)$ can be calculated by replacing $\Psi$ and $\Pi$
with $\psi$ and $\pi$ in Eqs.~\eqref{eq:7.5.37} and \eqref{eq:7.5.38}. This gives the equation
of motion
\begin{equation}
\dot\psi=\frac{\delta H_0}{\delta\pi}
=\gamma^0(\boldsymbol\gamma\cdot\boldsymbol\nabla+m)\psi
\tag{7.5.40}\label{eq:7.5.40}
\end{equation}
or more neatly
\begin{equation}
(\gamma^\mu\partial_\mu+m)\psi=0.
\tag{7.5.41}\label{eq:7.5.41}
\end{equation}
(The other equation of motion, $\dot\pi=-\delta H_0/\delta\psi$, yields just
the adjoint of this one.) Any field satisfying Eq.~\eqref{eq:7.5.41} can be written as
a Fourier transform
\begin{equation}
\psi(x)=(2\pi)^{-3/2}\int \dd^3p\sum_\sigma
\left\{u(\mathbf p,\sigma)e^{\ii p\cdot x}a_{\mathbf p,\sigma}
+v(\mathbf p,\sigma)e^{-\ii p\cdot x}b^\dagger_{\mathbf p,\sigma}\right\},
\tag{7.5.42}\label{eq:7.5.42}
\end{equation}
where $p^0=\sqrt{\mathbf p^2+m^2}$; $a_{\mathbf p,\sigma}$ and
$b^\dagger_{\mathbf p,\sigma}$ are operator coefficients; and
$u(\mathbf p,\mathord\pm\frac12)$ are the two independent solutions of
\begin{equation}
(\ii\gamma^\mu p_\mu+m)u(\mathbf p,\sigma)=0
\tag{7.5.43}\label{eq:7.5.43}
\end{equation}
and likewise
\begin{equation}
(-\ii\gamma^\mu p_\mu+m)v(\mathbf p,\sigma)=0
\tag{7.5.44}\label{eq:7.5.44}
\end{equation}
normalized so that\footnote{The matrix $\ii\gamma^\mu p_\mu$ has eigenvalues
$\pm m$, so $\sum u\bar u$ and $\sum v\bar v$ must be proportional to the
projection matrices $(-\ii\gamma^\mu p_\mu+m)/2m$ and
$(\ii\gamma^\mu p_\mu+m)/2m$, respectively. The proportionality factor may be
adjusted up to a sign by absorbing it in the definition of $u$ and $v$. The
overall sign is determined by positivity; $\operatorname{Tr}\sum u\bar u\beta
=\sum u^\dagger u$ and $\operatorname{Tr}\sum v\bar v\beta
=\sum v^\dagger v$ must be positive.}
\begin{equation}
\sum_\sigma u(\mathbf p,\sigma)\bar u(\mathbf p,\sigma)
=\frac{-\ii\gamma^\mu p_\mu+m}{2p^0},
\tag{7.5.45}\label{eq:7.5.45}
\end{equation}
\begin{equation}
\sum_\sigma v(\mathbf p,\sigma)\bar v(\mathbf p,\sigma)
=-\frac{\ii\gamma^\mu p_\mu+m}{2p^0}.
\tag{7.5.46}\label{eq:7.5.46}
\end{equation}
In order to obtain the desired anticommutators
\begin{align}
\{\psi_\alpha(\mathbf x,t),\bar\psi_\beta(\mathbf y,t)\}
&=\{\psi_\alpha(\mathbf x,t),\pi_\gamma(\mathbf y,t)\}
(\gamma^0)_{\gamma\beta}\notag\\
&=\ii(\gamma^0)_{\alpha\beta}\delta^3(\mathbf x-\mathbf y),
\tag{7.5.47}\label{eq:7.5.47}
\end{align}
\begin{equation}
\{\psi_\alpha(\mathbf x,t),\psi_\beta(\mathbf y,t)\}=0,
\tag{7.5.48}\label{eq:7.5.48}
\end{equation}
